\begin{prerequis}
Savoir distinguer une phrase affirmative d'une question; connaître les symboles
$=$, $\ne$, $<$, $\leq$; savoir produire un contre-exemple simple.
\end{prerequis}
\begin{automatismes}
Dire si chaque phrase est vraie, fausse, ou si elle ne possède pas de valeur de
vérité : « $7$ est premier »; « $x+2=5$ »; « calculer $3^2$ »; « tout carré est
positif »; « il existe un entier pair supérieur à $100$ ».
\end{automatismes}
\begin{activite}[Activité d'entrée -- Le tribunal des affirmations]
Par groupes, classer dix phrases proposées par le professeur. Pour chaque phrase
classée vraie ou fausse, rédiger une justification; pour chaque phrase refusée,
expliquer quelle information manque. Transformer ensuite deux phrases ouvertes
en propositions à l'aide de « pour tout » ou « il existe ».
\end{activite}

